\chapter{L'apparato digerente}

% ---------------------------------------------------------------------------
\section{Il peritoneo e la cavità addominale}

\textbf{Peritoneo}: membrana della cavità addominale, due strati continui $\rightarrow$
\textbf{viscerale} (avvolge gli organi) e \textbf{parietale} (riveste la parete). Organi per
rapporto col peritoneo:
\begin{itemize}
  \item \textbf{intraperitoneali}: stomaco, fegato, ileo (avvolti dal viscerale);
  \item \textbf{retroperitoneali}: reni, ureteri, aorta addominale (dietro il parietale);
  \item \textbf{retroperitoneali secondari}: pancreas e parte del duodeno (perdono il rivestimento
    nello sviluppo e si addossano alla parete posteriore).
\end{itemize}

% ---------------------------------------------------------------------------
\section{Generalità e funzioni}

Apparato digerente: processa il cibo, estrae i nutrienti, elimina i residui. Sei funzioni:
\begin{enumerate}
  \item \textbf{ingestione}: assunzione selettiva di cibo;
  \item \textbf{digestione}: frammentazione meccanica e chimica;
  \item \textbf{secrezione}: acidi, enzimi, tamponi (parete o organi accessori);
  \item \textbf{assorbimento}: molecole nutrienti nel sangue e nella linfa;
  \item \textbf{compattazione}: disidratazione del materiale non digeribile;
  \item \textbf{defecazione}: eliminazione delle feci.
\end{enumerate}
Il rivestimento difende da acidi/enzimi, stress meccanici e patogeni.

\rubrica{Organi e funzione primaria}
\begin{itemize}
  \item \textbf{cavità orale} (denti, lingua): trattamento meccanico, umidificazione, mescolamento
    con la saliva;
  \item \textbf{ghiandole salivari}: fluido lubrificante con enzimi che scindono i carboidrati;
  \item \textbf{faringe}: spinge il materiale nell'esofago; \textbf{esofago}: trasporto allo stomaco;
  \item \textbf{stomaco}: scissione chimica (acido, enzimi) e meccanica (contrazioni);
  \item \textbf{pancreas}: esocrino (tamponi ed enzimi) ed endocrino (ormoni);
  \item \textbf{fegato}: bile (digestione dei lipidi), deposito di nutrienti; \textbf{cistifellea}:
    deposito e concentrazione della bile;
  \item \textbf{intestino tenue}: digestione enzimatica e assorbimento; \textbf{crasso}:
    disidratazione e consolidamento dei residui.
\end{itemize}

% ---------------------------------------------------------------------------
\section{La cavità orale}

Prima porzione del canale alimentare, nello splancnocranio; comunica con l'esterno tramite la rima
buccale. Posteriormente $\rightarrow$ istmo delle fauci $\rightarrow$ faringe:
\begin{itemize}
  \item \textbf{rinofaringe}: comunica con le cavità nasali (coane);
  \item \textbf{orofaringe}: in continuità con la cavità orale;
  \item \textbf{laringofaringe}: continua con laringe ed esofago.
\end{itemize}
Comprende denti, gengive, lingua, palato, ghiandole. Due arcate gengivodentali (mascellare superiore,
mandibolare inferiore); vestibolo a ferro di cavallo.

\subsection{La lingua}
Organo muscolomembranoso; parte del pavimento orale. Funzioni: trattamento meccanico del cibo,
preparazione del bolo, analisi sensoriale (tattile, termica, gustativa), secrezione di enzimi e
mucine.
\begin{itemize}
  \item parti: radice, corpo, apice; solco mediano longitudinale; solco terminale a V (corpo/radice)
    col foro cieco;
  \item \textbf{papille}: filiformi e fungiformi (diffuse), foliate e vallate (presso il solco
    terminale);
  \item radice: \textbf{tonsilla linguale} (follicoli); tra radice ed epiglottide le pieghe
    glosso-epiglottiche $\rightarrow$ vallecole; archi palato-glossi/palato-faringei (tonsilla
    palatina).
\end{itemize}

\subsection{I denti}
Due ordini:
\begin{itemize}
  \item \textbf{decidui} (da latte): spuntano tra 6 e 30 mesi (dagli incisivi);
  \item \textbf{permanenti} (32): tra 6 e 25 anni; numerati per emiarcata (superiore 1--16 da dx a
    sx, inferiore 17--32 da sx a dx); tipi: incisivi, canini, premolari, molari.
\end{itemize}
Struttura:
\begin{itemize}
  \item parti: \textbf{corona} (sporgente), \textbf{radice} (1--3, nell'alveolo), \textbf{colletto}
    (intermedio); cavità pulpare con polpa (connettivo vascolarizzato e innervato);
  \item \textbf{parodonto}: gengiva + periostio alveolare + legamento connettivale (periostio--cemento);
  \item tessuti: \textbf{smalto} (copre la corona, calcificato, sali di fluoro, di origine
    epiteliale), \textbf{dentina} (massa del dente, odontoblasti con prolungamenti nei canalicoli),
    \textbf{cemento} (ricopre la radice, tessuto osseo con cementociti, più numerosi all'apice).
\end{itemize}

\subsection{Il palato}
\begin{itemize}
  \item \textbf{osseo/duro} (anteriore): processi palatini delle mascelle + lamine orizzontali dei
    palatini, con mucosa;
  \item \textbf{molle/velo palatino} (posteriore): lamina connettivale + muscoli (tensore ed elevatore
    del velo, palatoglosso, palato-faringeo, muscolo dell'ugola) + mucosa; centralmente l'\textbf{ugola
    palatina} (velo pendulo).
\end{itemize}

% ---------------------------------------------------------------------------
\section{Le ghiandole salivari}

Tre paia maggiori: \textbf{parotide}, \textbf{sottomandibolare}, \textbf{sottolinguale}; + numerose
minori. Producono $\sim$1--1,5 l/die. Secrezione sotto controllo del sistema nervoso autonomo (doppia
innervazione); riflesso salivare via trigemino e calici gustativi (nervi VII, IX, X).

\subsection{Parotide}
La più grande ($\sim$30 g), lobulata, con capsula fibrosa e tessuto adiposo; sotto la pelle, davanti
al lobo dell'orecchio (loggia parotidea, davanti allo sternocleidomastoideo). Secreto sieroso denso
con \textbf{amilasi salivare} (inizia la digestione dei carboidrati).

\subsection{Sottomandibolare}
A metà del corpo della mandibola, sotto il miloioideo. Secrezione mista sieromucosa a prevalenza
sierosa; adenomeri tubulari con componente mioepiteliale; \textbf{semilune sierose} (facilitano il
drenaggio della componente mucosa).

\subsection{Sottolinguale}
La più piccola (3--6 g), ovoidale, nel pavimento della bocca ($\sim$15 lobuli). In rapporto con
genioglosso, genioioideo, ventre anteriore del digastrico, miloioideo, pterigoideo mediale, osso
ioide. Dotti che sboccano nel solco sublinguale (presso la caruncola) o nel dotto sublinguale maggiore
(di Bartolini); vicino sbocca anche il dotto della sottomandibolare.

\subsection{La saliva}
99\% acqua + 1\% sostanze: elettroliti (sodio, potassio, magnesio, calcio), muco, sostanze
antibatteriche (IgA), ormoni (pepsina, testosterone), enzimi (alfa-amilasi, lipasi), cellule, batteri,
virus.

% ---------------------------------------------------------------------------
\section{Epiglottide, faringe e deglutizione}

\textbf{Epiglottide}: membrana cartilaginea che chiude le vie aeree nella deglutizione. Muscoli della
deglutizione:
\begin{itemize}
  \item \textbf{costrittori della faringe}: spingono il bolo verso l'esofago;
  \item \textbf{palato-faringeo e stilo-faringeo}: innalzano la laringe;
  \item \textbf{palatali}: innalzano il palato molle e la parete faringea.
\end{itemize}

\subsection{Le fasi della deglutizione}
\begin{enumerate}
  \item \textbf{buccale}: il bolo è spinto contro il palato duro; la retrazione della lingua lo spinge
    in faringe e innalza il palato molle; volontaria, poi (in orofaringe) parte il riflesso
    involontario;
  \item \textbf{faringea}: al contatto con archi palatali/parete posteriore, la laringe si innalza e
    l'epiglottide si ripiega $\rightarrow$ bolo oltre la glottide (chiusa) $\rightarrow$ i costrittori
    lo spingono nell'esofago;
  \item \textbf{esofagea}: apertura dello sfintere esofageo superiore $\rightarrow$ onde peristaltiche
    $\rightarrow$ apertura dello sfintere esofageo inferiore $\rightarrow$ stomaco.
\end{enumerate}

% ---------------------------------------------------------------------------
\section{L'esofago}

Tubo che collega faringe e stomaco (peristalsi); muscolatura controllata da riflessi viscerali. Organo
cavo, da C6 a T10, $\sim$25 cm; attraversa il mediastino posteriore, supera il diaframma, si apre nello
stomaco (T10).
\begin{itemize}
  \item segue la cifosi toracica; spinto a destra e poi in avanti per i rapporti con l'arco dell'aorta
    (a sinistra) e l'aorta toracica (posteriore);
  \item restringimenti: cervicale iniziale, toracico, orifizio esofageo del diaframma (transito
    rallentabile);
  \item drenaggio linfatico $\rightarrow$ linfonodi tracheali, mediastinici, tracheo-bronchiali,
    diaframmatici, gastrici di sinistra, celiaci $\rightarrow$ condotto toracico e giugulari interni.
\end{itemize}

\subsection{Nota clinica: reflusso gastro-esofageo}
Reflusso (acido) $\rightarrow$ la mucosa esofagea è esposta ad acido ed enzimi gastrici. Cause:
\begin{itemize}
  \item rilasciamento dello \textbf{sfintere esofageo inferiore} (SEI): favorito da fumo, caffè,
    farmaci o ernia iatale;
  \item rallentato svuotamento gastrico: pasti abbondanti, cibi grassi, cioccolato.
\end{itemize}

% ---------------------------------------------------------------------------
\section{Lo stomaco}

Organo cavo intraperitoneale, nello spazio sopra-mesocolico (ipocondrio sinistro ed epigastrio); a
sacco, schiacciato in senso anteroposteriore.
\begin{itemize}
  \item rapporti: avanti diaframma, fegato, parete addominale; dietro pancreas e milza (tra stomaco
    e pancreas la \textbf{borsa omentale});
  \item margini curvi = piccola e grande curvatura;
  \item irrorazione: rami del tripode celiaco (arterie gastriche e gastroepiploiche) e arteria
    splenica.
\end{itemize}

\subsection{Le porzioni}
\begin{itemize}
  \item \textbf{cardias}: attorno allo sbocco dell'esofago;
  \item \textbf{fondo}: a cupola, la parte più alta (concavità diaframmatica);
  \item \textbf{corpo}: la porzione più ampia, verticale;
  \item \textbf{parte pilorica} (antro + canale pilorico): conica, obliqua verso destra.
\end{itemize}
Funzioni: accumulo del cibo, trasformazione meccanica, digestione chimica (acidi ed enzimi).

\subsection{Parete e ghiandole gastriche}
\begin{itemize}
  \item mucosa: epitelio di cellule colonnari a secrezione mucosa + \textit{brush cell} (con
    microvilli, analizzano il contenuto); secreto mucoso alcalino difende dall'acido;
  \item ghiandole gastriche (corpo e fondo): cellule del colletto (mucose), principali, parietali,
    staminali, neuroendocrine;
  \item mucosa in pliche (lisce a fondo/cardias, rilevate ad antro/piloro) col plesso sottomucoso;
  \item tonaca muscolare a 3 strati (obliquo profondo, circolare intermedio, longitudinale
    superficiale) col plesso mienterico.
\end{itemize}

\subsection{Cellule delle ghiandole gastriche}
\begin{itemize}
  \item \textbf{mucose}: secernono muco;
  \item \textbf{parietali}: acido cloridrico ($\downarrow$ pH, uccide i microrganismi) + fattore
    intrinseco (assorbimento della vitamina B12) + grelina (appetito);
  \item \textbf{principali}: pepsinogeno (precursore della pepsina, digerisce le proteine) + lipasi
    gastrica (grassi del latte);
  \item \textbf{enterocromaffini}: ormoni digestivi $\rightarrow$ cellule G = gastrina (stimola la
    secrezione acida/enzimatica e la motilità).
\end{itemize}

% ---------------------------------------------------------------------------
\section{Il pancreas}

Ghiandola tubulo-acinosa composta, extramurale, all'altezza delle prime due vertebre lombari; a
martello (testa nella ``C'' duodenale, collo/istmo, corpo, coda verso l'alto e a sinistra).
\begin{itemize}
  \item \textbf{esocrina}: ghiandola tubuloacinosa a secrezione sierosa (simile alla parotide)
    $\rightarrow$ succo pancreatico;
  \item \textbf{endocrina}: \textbf{isole pancreatiche} $\rightarrow$ ormoni (insulina, glucagone)
    nel sangue, regolano il metabolismo dei carboidrati.
\end{itemize}

\subsection{Istologia e secrezione}
\textbf{Acino pancreatico} (unità esocrina): cellule acinose (enzimi digestivi, RER sviluppato, Golgi,
granuli di zimogeno) + cellule centroacinose. Secrezione: canalicoli $\rightarrow$ dotti intercalari
$\rightarrow$ intralobulari $\rightarrow$ dotto pancreatico principale $\rightarrow$ dotto coledoco.

\textbf{Succo pancreatico} (nel duodeno): acqua, ioni, enzimi:
\begin{itemize}
  \item \textbf{lipasi} (lipidi), \textbf{carboidrasi} (zuccheri e amidi), \textbf{nucleasi} (acidi
    nucleici);
  \item \textbf{proteolitici}: proteasi (grosse proteine) e peptidasi (peptidi $\rightarrow$
    amminoacidi).
\end{itemize}
Prodotto in fase cefalica, $\sim$1 l/die, pH 7,5--8,8 (ricco di bicarbonato dei dotti escretori).

% ---------------------------------------------------------------------------
\section{Il fegato}

Organo parenchimatoso, spazio sopramesocolico (ipocondrio destro, epigastrio, parte dell'ipocondrio
sinistro); intraperitoneale (quasi del tutto rivestito), ovoide.
\begin{itemize}
  \item vascolarizzazione all'ilo (faccia inferiore): arteriosa (arteria epatica propria) + venosa
    (vena porta);
  \item parenchima in \textbf{lobuli epatici} (aree poligonali, epatociti in cordoni);
  \item legamenti: falciforme e rotondo, piccolo omento (epatoduodenale, epatogastrico), gastrocolico;
  \item la ghiandola più voluminosa: $\sim$1,5 kg (+ $\sim$500 g di sangue) $\approx$ 2\% del peso
    corporeo.
\end{itemize}

\subsection{Funzioni}
\rubrica{Digestive e metaboliche}
\begin{itemize}
  \item sintesi di somatomedine; sintesi e secrezione di bile;
  \item accumulo di glicogeno e lipidi; controllo di glucosio, amminoacidi, acidi grassi ematici;
  \item interconversione di nutrienti (transaminazione, carboidrati $\rightarrow$ lipidi);
  \item sintesi e rilascio di colesterolo (con proteine vettrici);
  \item inattivazione di tossine; accumulo di ferro e vitamine liposolubili.
\end{itemize}
\rubrica{Altre funzioni}
\begin{itemize}
  \item sintesi di proteine plasmatiche e fattori della coagulazione;
  \item sintesi dell'angiotensinogeno (inattivo);
  \item fagocitosi degli eritrociti danneggiati (cellule di Kupffer); accumulo di sangue;
  \item degradazione di ormoni (insulina, adrenalina), immunoglobuline e farmaci liposolubili.
\end{itemize}
\rubrica{Regolazioni}
\begin{itemize}
  \item \textbf{metabolica}: tutto il sangue refluo dall'apparato digerente entra nel circolo portale
    $\rightarrow$ gli epatociti controllano i metaboliti;
  \item \textbf{ematologica}: principale riserva di sangue ($\sim$25\% della gittata cardiaca); i
    fagociti rimuovono eritrociti vecchi/patogeni/detriti, gli epatociti sintetizzano proteine
    plasmatiche;
  \item \textbf{bile}: sintetizzata dagli epatociti $\rightarrow$ accumulata nella cistifellea
    $\rightarrow$ duodeno; acqua e ioni (neutralizzano il chimo) + sali biliari (emulsionano i lipidi).
\end{itemize}

% ---------------------------------------------------------------------------
\section{La cistifellea}

Piccolo organo cavo piriforme (fino a 50--60 cm\textsuperscript{3} di bile), nella fossa cistica della
faccia viscerale del fegato; il collo continua nel dotto cistico.

Funzione: deposito e \textbf{concentrazione} della bile (concentrazione dei sali biliari). La bile
($\sim$1 l/die) $\rightarrow$ cistifellea $\rightarrow$ duodeno, soprattutto dopo l'arrivo di grassi,
stimolato dalla \textbf{colecistochinina (CCK)} che:
\begin{itemize}
  \item contrae la muscolatura della cistifellea e rilascia lo sfintere dell'ampolla;
  \item stimola il pancreas a secernere gli enzimi nel duodeno.
\end{itemize}
Bile (emulsione) + succhi pancreatici/enterici (idrolisi) $\rightarrow$ completano la digestione.
L'eccessiva concentrazione dei sali biliari $\rightarrow$ calcoli (colelitiasi).

% ---------------------------------------------------------------------------
\section{L'intestino tenue}

Assorbimento e digestione dei nutrienti (90\% nel tenue, 10\% nel crasso prossimale). Porzioni:
\begin{itemize}
  \item \textbf{duodeno}: breve ($\sim$25 cm), poco mobile, a ``C'';
  \item \textbf{tenue mesenteriale}: lungo ($\sim$6,5 m), mobile $\rightarrow$ digiuno e ileo (senza
    limite preciso).
\end{itemize}
Il peritoneo riveste il tenue mesenteriale e solo la parte superiore del duodeno (la discendente solo
nei 2/3 anteriori) $\rightarrow$ diversa mobilità. Il tenue mesenteriale è ripiegato in anse,
incorniciato dal colon; l'ileo confluisce nel cieco (valvola ileocecale). Funzione: degradazione del
chimo (carboidrati, proteine, lipidi) poi riassorbimento.

% ---------------------------------------------------------------------------
\section{L'intestino crasso}

Dalla valvola ileocecale all'ano; $\sim$1,6 m, prevalentemente addominale (tranne il retto, pelvico).
\begin{itemize}
  \item al cieco è annessa l'\textbf{appendice vermiforme};
  \item colon: ascendente, trasverso, discendente, sigmoideo (ascendente, discendente e retto solo
    parzialmente peritoneali);
  \item \textbf{retto}: ampolla rettale (pelvi) + canale anale (perineo) $\rightarrow$ ano.
\end{itemize}
Funzioni:
\begin{itemize}
  \item riassorbimento di acqua ed elettroliti $\rightarrow$ compattazione in feci;
  \item assorbimento di vitamine dalla flora batterica commensale;
  \item accumulo delle feci prima della defecazione.
\end{itemize}

\subsection{Nota clinica: patologie del colon-retto}
\begin{itemize}
  \item \textbf{polipi}: neoformazioni carnose della mucosa, benigne o maligne;
  \item \textbf{adenocarcinoma del colon-retto}: tumore maligno da proliferazione incontrollata delle
    cellule della mucosa.
\end{itemize}
